\documentclass[aspectratio=169]{beamer}
\usepackage[utf8]{inputenc}
\usepackage[english,russian]{babel}
\usepackage{cancel}
\usepackage{amssymb}
\usepackage{stmaryrd}
\usepackage{cmll}
\usepackage{graphicx}
\usepackage{amsthm}
\usepackage{tikz}
\usepackage{multicol}
\usetikzlibrary{patterns}
\usepackage{chronosys}
\usepackage{proof}
\usepackage{multirow}
\usepackage{comment}
\setbeamertemplate{navigation symbols}{}
%\usetheme{Warsaw}

\newtheorem{thm}{Теорема}[section]
\newtheorem{dfn}{Определение}[section]
\newtheorem{lmm}{Лемма}[section]
\newtheorem{exm}{Пример}[section]
\newtheorem{snote}{Пояснение}[section]

\newcommand{\divisible}%
{\mathrel{\lower.2ex%
\vbox{\baselineskip=0.7ex\lineskiplimit=0pt%
\kern6pt \hbox{.}\hbox{.}\hbox{.}}%
}}

\begin{document}

\newcommand\doubleplus{+\kern-1.3ex+\kern0.8ex}
\newcommand\mdoubleplus{\ensuremath{\mathbin{+\mkern-10mu+}}}

\begin{comment}
\begin{frame}{Уточнение исчисления предикатов}
\begin{itemize}
\item Пусть требуется доказывать утверждения про равенство. Введём $E(p,q)$ --- предикат <<равенство>>.\pause
\item Однако $\not\vdash E(p,q)\to E(q,p)$: если $D = \{0,1\}$ и $E(p,q) ::= (p>q)$,
то $\not\models E(p,q)\to E(q,p)$.\pause
\item Конечно, можем указывать $\forall p.\forall q.E(p,q)\to E(q,p) \vdash \varphi$.\pause
\item Но лучше добавим аксиому $\forall p.\forall q.E(p,q)\to E(q,p)$.\pause
\item Добавив необходимые аксиомы, получим \emph{теорию первого порядка}.
\end{itemize}
\end{frame}
\end{comment}

\begin{frame}{Теория первого порядка}
\begin{dfn}
Теорией первого порядка назовём исчисление предикатов с дополнительными (<<нелогическими>>
или <<математическими>>):
\begin{itemize}
\item предикатными и функциональными символами;
\item аксиомами.
\end{itemize}

Сущности, взятые из исходного исчисления предикатов, назовём \emph{логическими}
\end{dfn}
\end{frame}

\begin{frame}{Порядок логики/теории}
\begin{tabular}{llll}
Порядок & Кванторы & Формализует суждения\dots & Пример\\\hline
нулевой & запрещены & об отдельных значениях & И.В.\\
первый & по предметным переменным & о множествах & И.П.\\
%    &   & $S = \{ t\ |\ \psi[x := t] \}$ \\
    &   \multicolumn{2}{l}{\color{olive}$\{2,3,5,7,\dots\} = \{ t\ |\ \forall p.\forall q.(p \ne 1 \with q \ne 1) \rightarrow (t \ne p\cdot q)\}$}\\
второй & по предикатным переменным & о множествах множеств & Типы\\
    &   \multicolumn{2}{l}{\color{olive}$S = \{ \{t\ |\ P(t)\}\ |\ \varphi[p := P] \}$}\\
 & \dots 
\end{tabular}
\pause\vspace{0.3cm}

\begin{exm}[логики 2 порядка]\upshape
\begin{tabular}{ll}
$\alpha\rightarrow\beta\rightarrow\alpha$ (сх. акс. 1) & $\forall a.\forall b.a \rightarrow b \rightarrow a$ \vspace{0.1cm}\\
\texttt{let rec map f l = match l with} & $map: \forall a.\forall b.(a \rightarrow b) \rightarrow a\texttt{ list} \rightarrow b\texttt{ list}$ \\
\texttt{| [] -> []} \\
\texttt{| l1::ls -> f l1 :: map f ls}\vspace{0.1cm}\\
\texttt{map ((+) 1) [1;2;3] = [2;3;4]}
\end{tabular}
\end{exm}

\end{frame}

\begin{frame}{Формальная арифметика}
\begin{dfn}
Формальная арифметика --- теория первого порядка, со следующими добавленными нелогическими \dots
\begin{itemize}
\item двухместными функциональными символами $(+)$, $(\cdot)$; одноместным функциональным символом $(')$, 
нульместным функциональным символом $0$;\pause
\item двухместным предикатным символом $(=)$;\pause
\item восемью нелогическими \emph{аксиомами}:\vspace{0.1cm}
\begin{tabular}{ll}
(A1) $a=b \to a=c \to b=c$             &(A5) $a+0 = a$                     \\
(A2) $a=b \to a'=b'$                   &(A6) $a+b' = (a+b)'$               \\
(A3) $a'=b' \to a=b$                   &(A7) $a\cdot 0 = 0$                \\
(A4) $\neg a' = 0$                     &(A8) $a\cdot b' = a \cdot b + a$
\end{tabular}\pause
\item нелогической схемой аксиом индукции $\psi[x:=0]\with(\forall x.\psi\to \psi[x:=x'])\to \psi$ с метапеременными $x$ и $\psi$.
\end{itemize}
\end{dfn}
\end{frame}

\begin{comment}
\begin{frame}{Докажем, что $a=a$}
\small
Пусть $\top ::= 0=0\to 0=0 \to 0=0$, тогда:\pause

\begin{tabular}{lll}
(1) & $a=b\to a=c \to b=c$ & (Акс. А1)\\
(2) & $(a=b\to a=c \to b=c) \to \top \to (a=b\to a=c \to b=c)$ & (Сх. акс. 1)\\
(3) & $\top \to (a=b\to a=c \to b=c)$ & (M.P. 1, 2)\\
(4) & $\top \to (\forall c.a = b\to a = c \to b = c)$ & (Введ. $\forall$)\\\pause
(5) & $\top \to (\forall b.\forall c.a = b\to a = c \to b = c)$ & (Введ. $\forall$)\\
(6) & $\top \to (\forall a.\forall b.\forall c.a = b\to a = c \to b = c)$ & (Введ. $\forall$)\\\pause
(7) & $\top$ & (Сх. акс 1)\\
(8) & $(\forall a.\forall b.\forall c.a = b\to a = c \to b = c)$ & (M.P. 7, 6)\\\pause
(9) & $(\forall a.\forall b.\forall c.a = b\to a = c \to b = c) \to $\\
    & $\to (\forall b.\forall c.a+0 = b\to a+0 = c \to b = c)$ & (Сх. акс. 11)\\\pause
(10) & $\forall b.\forall c.a+0 = b\to a+0 = c \to b = c$ & (M.P. 8, 9)\\\pause
(12) & $\forall c.a+0 = a\to a+0 = c \to a = c$ & (M.P. 10, 11)\\\pause
(14) & $a+0 = a\to a+0 = a \to a = a$ & (M.P. 12, 13)\\\pause
(15) & $a+0 = a$ & (Акс. А5)\\
(16) & $a+0 = a \to a = a$ & (M.P. 15, 14)\\
(17) & $a = a$ & (M.P. 15, 16)
\end{tabular}
\end{frame}
\end{comment}


\begin{frame}{Выразимость отношений в Ф.А.}
\begin{dfn}
Будем говорить, что отношение $R\subseteq \mathbb{N}^n_0$ выразимо в ФА, 
если существует формула $\rho$, что:
\begin{enumerate}
\item если $\langle a_1,\dots,a_n \rangle \in R$, то $\vdash \rho(\overline{a_1},\dots,\overline{a_n})$
\item если $\langle a_1,\dots,a_n \rangle \notin R$, то $\vdash \neg\rho(\overline{a_1},\dots,\overline{a_n})$
\end{enumerate}
\end{dfn}\pause

\begin{thm} отношение <<равно>> выразимо в Ф.А.:
$R = \{ \langle x,x \rangle\ |\ x \in \mathbb{N}_0 \}$
\end{thm}\pause

\begin{proof}
Пусть $\rho := x_1=x_2$. \pause Тогда:
\begin{itemize}
\item $\vdash p = p$ при $p := \overline{k}$ при всех $k \in \mathbb{N}_0$: \pause $\vdash 0=0$, $\vdash 0'=0'$, $\vdash 0''=0''$, ... \pause
\item $\vdash \neg p = q$ при $p := \overline{k}$, $q := \overline{s}$ при всех $k,s \in \mathbb{N}_0$ и $k \ne s$.

\pause $\vdash \neg 0 = 0'$, $\vdash \neg 0 = 0''$, $\vdash \neg 0''' = 0'$, ...
\end{itemize}
\end{proof}

\end{frame}

\begin{frame}{Представимость функций в Ф.А.}
\begin{dfn}
Будем говорить, что функция $f: \mathbb{N}^n_0\to\mathbb{N}_0$ представима в ФА, 
если существует формула $\varphi$, что:
\begin{enumerate}
\item если $f(a_1,\dots,a_n) = u$, то $\vdash \varphi(\overline{a_1},\dots,\overline{a_n},\overline{u})$
\item если $f(a_1,\dots,a_n) \ne u$, то $\vdash \neg\varphi(\overline{a_1},\dots,\overline{a_n},\overline{u})$
\item для всех $a_i \in \mathbb{N}_0$ выполнено $\vdash (\exists x.\varphi(\overline{a_1},\dots,\overline{a_n},x)) \with 
   (\forall p.\forall q.\varphi(\overline{a_1},\dots,\overline{a_n},p)\with \varphi(\overline{a_1},\dots,\overline{a_n},q)\rightarrow p=q)$
\end{enumerate}
\end{dfn}
\end{frame}

\begin{frame}{Соответствие рекурсивных и представимых функций}
\begin{thm}Любая рекурсивная функция представима в Ф.А.\end{thm}\pause

\begin{thm}Любая представимая в Ф.А. функция рекурсивна.\end{thm}
\end{frame}

\begin{frame}{Примитивы Z, N, U представимы в Ф.А.}
\begin{thm}
Примитивы $Z$, $N$ и $U^k_n$ представимы в Ф.А.
\end{thm}\pause

\begin{proof}
\begin{itemize}
\item $\zeta(x_1,x_2) := x_2=0$, \pause формальнее: $\zeta(x_1,x_2) := x_1=x_1 \with x_2=0$\pause
\item $\nu(x_1,x_2) := x_2=x_1'$\pause
\item $\upsilon(x_1,\dots,x_n,x_{n+1}) := x_k = x_{n+1}$ \pause \vspace{0.1cm}

формальнее: 
     $\upsilon(x_1,\dots,x_n,x_{n+1}) := (\underset{i\ne k,n+1}{\with} x_i=x_i) \with x_k = x_{n+1}$
\end{itemize}
\end{proof}
\end{frame}

\begin{frame}{Примитив S представим в Ф.А.}

$$S\langle f,g_1,\dots,g_k\rangle(x_1,\dots,x_n) = f(g_1(x_1,\dots,x_n),\dots,g_k(x_1,\dots,x_n))$$\pause
\begin{thm}
Пусть функции $f,g_1,\dots,g_k$ представимы в Ф.А. Тогда $S\langle f,g_1,\dots,g_k \rangle$ представима в Ф.А.
\end{thm}\pause

\begin{proof}
Пусть $f$, $g_1$, ..., $g_k$ представляются формулами $\varphi$, $\gamma_1$, ..., $\gamma_k$. \pause

Тогда 
$S\langle f,g_1,\dots,g_k\rangle$ будет представлена формулой
$$\exists g_1.\dots.\exists g_k.\varphi(g_1,\dots,g_k,x_{n+1})\with\gamma_1(x_1,\dots,x_n,g_1)\with\dots\with\gamma_k(x_1,\dots,x_n,g_k)$$
\end{proof}

\end{frame}

\begin{frame}{$\beta$-функция Гёделя}
Задача: закодировать последовательность натуральных чисел произвольной длины.\pause
\begin{dfn}$\beta$-функция Гёделя: $\mathcal{\beta}(b,c,i) := b \% (1 + (i+1) \cdot c)$\\
Здесь (\%) --- остаток от деления.
\end{dfn}\pause
\begin{thm}$\beta$-функция Гёделя представима в Ф.А. формулой
$$\hat{\beta}(b,c,i,d) := \exists q.(b = q \cdot (1 + c \cdot (i+1)) + d) \& (d < 1 + c \cdot (i+1))$$
\end{thm}\pause\vspace{-0.3cm}
Деление $b$ на $x$ с остатком: найдутся частное $(q)$ и остаток $(d)$, что
$b = q\cdot x + d$ и $0 \le d < x$.
\pause
\begin{thm}Если $a_0, \dots, a_n \in \mathbb{N}_0$, то найдутся такие $b,c \in \mathbb{N}_0$, что
$a_i = \beta(b,c,i)$
\end{thm}
\end{frame}

\begin{frame}{Доказательство свойства $\beta$-функции}
\begin{thm}Китайская теорема об остатках (вариант формулировки): если $u_0, \dots, u_n$ --- попарно 
взаимно просты, и $0 \le a_i < u_i$, то существует такой $b$, что $a_i = b \% u_i$.
\end{thm}\pause
Положим $c = \max(a_0,\dots,a_n,n)!$ и $u_i = 1+c\cdot(i+1)$.\pause
\begin{itemize}
\item $\text{НОД}(u_i,u_j) = 1$, если $i \ne j$.\pause

Пусть $p$ --- простое, $u_i\divisible p$ и $u_j\divisible p$ ($i < j$). \pause
Заметим, что $u_j-u_i = c \cdot (j-i)$. Значит, $c\divisible p$ или $(j-i)\divisible p$. \pause
Так как $j-i \le n$, то $c \divisible (j-i)$, потому если и $(j-i)\divisible p$, всё равно $c \divisible p$. \pause
Но и $(1+c\cdot(i+1))\divisible p$, отсюда $1 \divisible p$ --- что невозможно. \pause
\item $0 \le a_i < u_i$.\pause
\end{itemize}
Условия китайской теоремы об остатках выполнены и найдётся $b$, что \vspace{-0.2cm}
$$a_i = b \% (1 + c\cdot(i+1)) = \beta(b,c,i)$$
\qed
\end{frame}

\begin{frame}{Примитив <<примитивная рекурсия>> представим в Ф.А.}

Пусть $f: \mathbb{N}^n_0\to\mathbb{N}_0$ и $g:\mathbb{N}^{n+2}_0\to\mathbb{N}_0$ представлены формулами $\varphi$ и $\gamma$. 

Зафиксируем $x_1, \dots, x_n, y \in \mathbb{N}_0$.\pause\vspace{0.3cm}

\begin{tabular}{lll}
Шаг вычисления & Об. & Утверждение в Ф.А.\\
$R\langle f,g\rangle (x_1,\dots,x_n,0) = f(x_1,\dots,x_n)$ & $a_0$ & $\vdash \varphi(\overline{x_1},\dots,\overline{x_n},\overline{a_0})$\\ \pause
$R\langle f,g\rangle (x_1,\dots,x_n,1) = g(x_1,\dots,x_n,0,a_0)$ & $a_1$ & $\vdash \gamma(\overline{x_1},\dots,\overline{x_n},0,\overline{a_0},\overline{a_1})$\\ \pause
$\dots$\\
$R\langle f,g\rangle (x_1,\dots,x_n,y) = g(x_1,\dots,x_n,y-1,a_{y-1})$ & $a_y$ & $\vdash \gamma(\overline{x_1},\dots,\overline{x_n},\overline{y-1},\overline{a_{y-1}},\overline{a_y})$
\end{tabular}\pause\vspace{0.3cm}

По свойству $\beta$-функции, найдутся $b$ и $c$, что
$\beta (b,c,i) = a_i$ для $0 \le i \le y$.\pause

\begin{thm}Примитив $R\langle f,g\rangle$ представим в Ф.А. формулой $\rho(x_1,\dots,x_n,y,a)$:\vspace{-0.1cm}
$$\begin{array}{l}\exists b. \exists c. (\exists a_0. \hat{\beta} (b,c,0,a_0) \& \varphi (x_1,...x_n, a_0)) \\
       \&\;\;\;\;\forall k.k < y \rightarrow \exists d . \exists e . \hat{\beta} (b,c,k,d) \& \hat{\beta} (b,c,k',e) \& \gamma (x_1,..x_n,k,d,e) \\
       \&\;\;\;\;\hat{\beta} (b,c,y,a) 
\end{array}$$
\end{thm}
\end{frame}

\begin{frame}{Представимость рекурсивных функций в Ф.А.}
\begin{thm}Пусть функция $f:\mathbb{N}^{n+1}_0 \to \mathbb{N}_0$ представима в Ф.А.
формулой $\varphi(x_1,\dots,x_{n},y,r)$. Тогда примитив $M\langle f\rangle$ представим в Ф.А. 
формулой $$\mu(x_1,\dots,x_n,y) := \varphi(x_1,\dots,x_n,y,0) \with \forall u.u < y \to \neg\varphi(x_1,\dots,x_n,u,0)$$
\end{thm}\pause

\begin{thm}Если $f$ --- рекурсивная функция, то она представима
в Ф.А.
\end{thm}
\begin{proof}Индукция по структуре $f$.
\end{proof}
\end{frame}

\begin{frame}{Рекурсивность представимых функций в Ф.А.}
Фиксируем $f$ и $x_1, x_2, \dots, x_n$. Обозначим $y = f(x_1,x_2,\dots,x_n)$. \pause
По представимости нам известна $\varphi$, что $\vdash \varphi(\overline{x_1},\overline{x_2},\dots,\overline{x_n},\overline{y})$. \pause
Давайте просто переберём все результаты и доказательства!\pause

\begin{enumerate}
\item Закодируем доказательства натуральными числами.\pause
\item Напишем рекурсивную функцию, проверяющую доказательства на корректность.\pause
\item Параллельный перебор значений и доказательств: $s = 2^y \cdot 3^p$. Переберём все $s$, по $s$ получим $y$ и $p$.
Проверим, что $p$ --- код доказательства $\vdash \varphi(\overline{x_1},\overline{x_2},\dots,\overline{x_n},\overline{y})$.
\end{enumerate}
\end{frame}

\begin{frame}{Гёделева нумерация}

\begin{enumerate}
\item Отдельный символ.\vspace{0.2cm}

\begin{tabular}{lc|lc||cll}
Номер & Символ & Номер & Символ & Имя & $k,n$ & Гёделев номер\\
3 & ( &               17 & $\&$ &  0 & $0,0$ & $27 + 6$\\
5 & ) &               19 & $\forall$ & $(')$ & $0,1$ & $27 + 6 \cdot 3$\\
7 & , &               21 & $\exists$ & $(+)$ & $0,2$ & $27 + 6 \cdot 9$\\
9 & . &               23 & $\vdash$ & $(\cdot)$ & $1,2$ & $27 + 6 \cdot 2 \cdot 9$\\
11 & $\neg$ &         $25 + 6\cdot k$ & $x_k$ & $(=)$ & $0,2$ & $29 + 6 \cdot 9$\\
13 & $\rightarrow$ &  $27 + 6\cdot 2^k \cdot 3^n$ & $f_k^n$\\
15 & $\vee$ &         $29 + 6\cdot 2^k \cdot 3^n$ & $P_k^n$ 
\end{tabular}\pause

\item Формула. $\phi \equiv s_0s_1\dots s_{n-1}$. Гёделев номер: $\ulcorner\phi\urcorner = 2^{\ulcorner s_0\urcorner}\cdot 3^{\ulcorner s_1\urcorner} 
\cdot \dots \cdot p_{n-1}^{\ulcorner s_{n-1}\urcorner}$.\pause

\item Доказательство. $\Pi = \delta_0\delta_1\dots\delta_{k-1}$, его гёделев номер: $\ulcorner\Pi\urcorner =
2^{\ulcorner \delta_0\urcorner}\cdot 3^{\ulcorner \delta_1\urcorner} \cdot \dots \cdot p_{k-1}^{\ulcorner \delta_{k-1}\urcorner}$
\end{enumerate}
\end{frame}

\begin{frame}{Проверка доказательства на корректность}
\begin{thm}Следующая функция рекурсивна:
$$\text{proof}(f,x_1,x_2,\dots,x_n,y,p) = \left\{\begin{array}{ll} 
0, & \mbox{если} \vdash\phi(\overline{x_1},\overline{x_2},\dots,\overline{x_n},\overline{y}),\\
   & p \mbox{ --- гёделев номер вывода}, f = \ulcorner\phi\urcorner \\ 
1, & \mbox{иначе}
\end{array}\right.$$\vspace{-0.1cm}\pause
\end{thm}
\begin{proof}[Идея доказательства]
\begin{enumerate}
%\item По представимости, $\vdash\phi(\overline{x_1},\overline{x_2},\dots,\overline{x_n},\overline{y})$ 
%эквивалентно $f(x_1,x_2,\dots,x_n)=y$.\pause
\item Проверка доказательства вычислима.\pause
\item Согласно тезису Чёрча, любая вычислимая функция
вычислима с помощью рекурсивных функций.
\end{enumerate}
\end{proof}
\end{frame}

\begin{frame}{Перебор доказательств}
\begin{lmm}Следующие функции рекурсивны:
\begin{enumerate}
\item Функции $\text{plog}_k(n) = \max\{p: n \divisible k^p\}$, $\text{fst}(x) = \text{plog}_2(x)$ и $\text{snd}(x) = \text{plog}_3(x)$.\pause
\item Числовые литералы: $\overline{k}: \mathbb{N}_0 \rightarrow \mathbb{N}_0$, $\overline{k}(x) = k$.\pause
\end{enumerate}
\end{lmm}
\begin{thm}Если $f: \mathbb{N}^n_0\rightarrow\mathbb{N}_0$, и $f$ представима в Ф.А. формулой $\varphi$, то $f$ --- рекурсивна.
\end{thm}\pause
\begin{proof}
Пусть заданы $x_1,x_2,\dots,x_n$. Ищем $\langle y, p\rangle$, 
что $\text{proof}(\ulcorner\varphi\urcorner,x_1,x_2,\dots,x_n,y,p)=1$, \pause
напомним: $y = f(x_1,x_2,\dots,x_n)$, $p = \ulcorner\Pi\urcorner$,
$\Pi$ --- доказательство $\varphi(\overline{x_1},\overline{x_2},\dots,\overline{x_n},\overline{y})$. \pause
$$f = S \langle \text{fst}, M\langle S \langle \text{proof}, \overline{\ulcorner\varphi\urcorner}, U^1_{n+1}, U^2_{n+1}, \dots, U^n_{n+1}, 
  S \langle \text{fst},U^{n+1}_{n+1}\rangle, S\langle \text{snd}, U^{n+1}_{n+1}\rangle \rangle \rangle \rangle$$
\end{proof}
\end{frame}

\begin{frame}{}
\LARGE\begin{center}Теоремы Гёделя о неполноте арифметики\end{center}
\end{frame}

\begin{frame}{Основные свойства исчислений: Ф.А.}
\begin{tabular}{lllll}
      & К.И.В. & И.И.В. & К.И.П. & Ф.А. + кл. модель\\\hline
корректность & да  & да  & да (лекция 5) & да (сейчас)\\
непротиворечивость & да  & да  & да (лекция 6) & {\color{red}верим} (т. Гёделя №2)\\
полнота & да & да & да (лекция 6) & {\color{red}нет} (т. Гёделя №1)\\
разрешимость & да  & да & нет (лекция 7)& {\color{red}нет} (док-во т. Тарского)
\end{tabular}
\end{frame}

\begin{frame}{Классическая модель Ф.А.}
А как определять <<нестандартные>> предикаты и функции ($Q'_1$, $c(p,q)$ и т.п.)? \pause
Для простоты разрешим только нелогические функциональные и предикатные символы ($=$, $+$, $\cdot$, 0, $'$).\pause
\begin{dfn}Классическая модель формальной арифметики: $D = \mathbb{N}_0$, оценки символов --- в соответствии
с аксиоматикой Пеано.
\end{dfn}

\begin{thm}Формальная арифметика корректна\end{thm}
\begin{proof}Свойства аксиом $A1\dots A8$ очевидны. 

Доказательство схемы аксиом индукции:
$$\psi(0) \with (\forall x.\psi(x) \rightarrow \psi(x')) \rightarrow \psi(x)$$
Математическая индукция по $x$.
\end{proof}
\end{frame}

\begin{frame}{Схема аксиом индукции чуть подробнее}
Фиксируем некоторую формулу $\psi$, 
и рассмотрим схему аксиом $$\psi(0) \with (\forall x.\psi(x) \rightarrow \psi(x')) \rightarrow \psi(x)$$

Разберём случаи:
\begin{enumerate}
\item $\llbracket\psi(0)\rrbracket = \text{Л}$ либо 
$\llbracket\forall x.\psi(x) \rightarrow \psi(x')\rrbracket = \text{Л}$. Тогда антецедент импликации ложен 
и формула истинна.
\item Антецедент импликации истинен. Тогда докажем индукцией по $x$, что $\llbracket \psi(x) \rrbracket = \text{И}$.
База:
$x = 0$.

Переход: $x := s'$. Тогда $s'$ раз применяем переход $$\llbracket \theta_0(x) = \theta_1(x) \rightarrow \theta_0(x') = \theta_1(x') \rrbracket^{x := \overline{0\dots s}} = \text{И}$$
отсюда $$\llbracket\theta_0(x') = \theta_1(x') \rrbracket^{x := s} = \llbracket\theta_0(x) = \theta_1(x) \rrbracket^{x := s+1} = \text{И}$$
\end{enumerate}

\pause
Можно ли верить этому доказательству (доказываем индукцию через индукцию)?
\end{frame}

\begin{frame}{Самоприменимость}
\begin{dfn}Пусть $\xi$ --- формула с единственной свободной
переменной $x_1$. Тогда:
$\langle\ulcorner \xi \urcorner,p\rangle \in W_1$, если $\vdash \xi(\overline{\ulcorner \xi \urcorner})$ и $p$ --- номер доказательства.
\end{dfn}

\begin{dfn}Отношение $W_1$ рекурсивно, поэтому выражено в Ф.А. формулой $\omega_1$ со свободными переменными $x_1$ и $x_2$, причём:
\begin{enumerate}
\item $\vdash \omega_1(\overline{\ulcorner \varphi \urcorner},\overline{p})$, если $p$ --- гёделев номер
доказательства самоприменения $\varphi$;
\item $\vdash \neg\omega_1(\overline{\ulcorner \varphi \urcorner},\overline{p})$ иначе.
\end{enumerate}
\end{dfn}

\begin{dfn}
Определим формулу $\sigma(x_1) := \forall p.\neg\omega_1(x_1,p)$. 
\end{dfn}
\end{frame}

\begin{frame}{Первая теорема Гёделя о неполноте арифметики}
\begin{dfn}Если для любой формулы $\phi(x)$ из $\vdash\phi(0)$, $\vdash\phi(\overline{1})$,
$\vdash\phi(\overline{2})$, $\dots$ выполнено $\not\vdash\exists x.\neg\phi(x)$, 
то теория \emph{омега-непротиворечива}.
\end{dfn}

\begin{thm}{Первая теорема Гёделя о неполноте арифметики}
\begin{itemize}
\item Если формальная арифметика непротиворечива, то $\not\vdash\sigma(\overline{\ulcorner\sigma\urcorner})$.
\item Если формальная арифметика $\omega$-непротиворечива, то $\not\vdash\neg\sigma(\overline{\ulcorner\sigma\urcorner})$.
\end{itemize}
\end{thm}
\end{frame}

\begin{frame}{Доказательство теоремы Гёделя}
Напомним: $\sigma(x_1) := \forall p.\neg\omega_1(x_1,p)$. $W_1(\ulcorner\xi\urcorner,p)$ --- $p$ есть доказательство самоприменения $\xi$.
\begin{proof}
\begin{itemize}
\item Пусть $\vdash\sigma(\overline{\ulcorner\sigma\urcorner})$. Значит, $p$ --- номер доказательства. \pause Тогда
$\langle\ulcorner\sigma\urcorner,p\rangle \in W_1$. \pause Тогда $\vdash\omega_1(\overline{\ulcorner\sigma\urcorner},\overline{p})$. \pause 
Тогда $\vdash\exists p.\omega_1(\overline{\ulcorner\sigma\urcorner},p)$. \pause То есть
$\vdash\neg\forall p.\neg\omega_1(\overline{\ulcorner\sigma\urcorner},p)$. \pause То есть $\vdash\neg\sigma(\overline{\ulcorner\sigma\urcorner})$. Противоречие.
\pause
\item Пусть $\vdash\neg\sigma(\overline{\ulcorner\sigma\urcorner})$. \pause То есть $\vdash\exists p.\omega_1(\overline{\ulcorner\sigma\urcorner},p)$.
\begin{itemize}
\item
\pause Но найдётся ли натуральное число $p$, что $\vdash\omega_1(\overline{\ulcorner\sigma\urcorner},\overline{p})$?
\pause Пусть нет. То есть $\vdash\neg\omega_1(\overline{\ulcorner\sigma\urcorner},\overline{0})$,
                          $\vdash\neg\omega_1(\overline{\ulcorner\sigma\urcorner},\overline{1})$, 
                          \dots \pause По $\omega$-непротиворечивости $\not\vdash\exists p.\neg\neg\omega_1(\overline{\ulcorner\sigma\urcorner},p)$. \pause
\end{itemize}
       Значит, найдётся натуральное $p$, что $\vdash\omega_1(\overline{\ulcorner\sigma\urcorner},\overline{p})$. \pause 
То есть, $\langle\ulcorner\sigma\urcorner, p\rangle\in W_1$. \pause
То есть, $p$ --- доказательство самоприменения $\sigma$: $\vdash\sigma(\overline{\ulcorner\sigma\urcorner})$. \pause Противоречие.
\end{itemize}
\end{proof}
\end{frame}

\begin{frame}{Почему теорема о неполноте?}
\begin{dfn}\emph{Семантически} полная теория --- теория, в которой любая общезначимая формула доказуема.\\
\emph{Синтаксически} полная теория --- теория, в которой для каждой формулы $\alpha$ выполнено $\vdash\alpha$ или $\vdash\neg\alpha$.\end{dfn}
\begin{thm}Формальная арифметика с классической моделью семантически неполна.\end{thm} \pause
\begin{proof}
 Рассмотрим Ф.А. с классической моделью. \pause
 Из теоремы Гёделя имеем $\not\vdash\sigma(\overline{\ulcorner\sigma\urcorner})$. \pause
 Рассмотрим $\sigma(\overline{\ulcorner\sigma\urcorner}) \equiv \forall p.\neg\omega_1(\overline{\ulcorner\sigma\urcorner},p)$:
нет числа $p$, что $p$ --- номер доказательства $\sigma(\overline{\ulcorner\sigma\urcorner})$. \pause 
 То есть, $\llbracket \forall p.\neg\omega_1(\overline{\ulcorner\sigma\urcorner},p) \rrbracket = \text{И}$. \pause
 То есть, $\models \sigma(\overline{\ulcorner\sigma\urcorner})$.
\end{proof}

\end{frame}

\begin{frame}{Первая теорема Гёделя о неполноте в форме Россера}
\begin{dfn}$\theta_1\le\theta_2 \equiv (\exists p.p+\theta_1=\theta_2)\quad\quad\theta_1<\theta_2\equiv\theta_1\le\theta_2\with\neg\theta_1=\theta_2$\end{dfn}\pause
\begin{dfn}Пусть $\langle \ulcorner\xi\urcorner,p\rangle \in W_2$, если $\vdash\neg\xi(\overline{\ulcorner\xi\urcorner})$ и $p$ --- номер доказательства.
Пусть $\omega_2$ выражает $W_2$ в формальной арифметике.\end{dfn}\pause
\begin{thm}Рассмотрим $\rho(x_1) = \forall p.\omega_1(x_1,p)\rightarrow\exists q.q \le p \with \omega_2(x_1,q)$. \pause
Тогда $\not\vdash\rho(\overline{\ulcorner\rho\urcorner})$ и $\not\vdash\neg\rho(\overline{\ulcorner\rho\urcorner})$. \pause
$\rho(\overline{\ulcorner\rho\urcorner})$: <<Меня легче опровергнуть, чем доказать>>
\end{thm}%\pause
%\begin{proof}Сложное техническое доказательство, использующее утверждения
%вида $\vdash a \le \overline{n} \leftrightarrow a = \overline{0} \vee a = \overline{1} \vee \dots \vee a = \overline{n}$
%\end{proof} \pause
\end{frame}

\begin{frame}{Формальное доказательство}
Неполнота варианта теории, изложенной выше, формально доказана на Coq, Russell O'Connor, 2005:\\
\vspace{0.5cm}
``My proof, excluding standard libraries and the library for Pocklington’s criterion, 
consists of 46 source files, 7 036 lines of specifications, 37 906 lines of proof, 
and 1 267 747 total characters. The size of the gzipped tarball (gzip -9) of all 
the source files is 146 008 bytes, which is an estimate of the information content of my proof.''\\
\vspace{0.5cm}
\texttt{Theorem Incompleteness : forall T : System,}\\
\hspace{0.5cm}\texttt{Included Formula NN T ->}\\
\hspace{0.5cm}\texttt{RepresentsInSelf T ->}\\
\hspace{0.5cm}\texttt{DecidableSet Formula T ->}\\
\hspace{0.5cm}\texttt{exists f : Formula,}\\
\hspace{0.5cm}\texttt{Sentence f/\textbackslash (SysPrf T f \textbackslash/ SysPrf T (notH f) -> Inconsistent LNN T)}.
\end{frame}

\begin{frame}{Consis}
\begin{lmm}
$\vdash 1=0$ тогда и только тогда, когда $\vdash\alpha$ при любом $\alpha$.
\end{lmm}
\begin{dfn}
Обозначим за $\psi(x,p)$ формулу, выражающую в формальной арифметике рекурсивное
отношение Proof: $\langle \ulcorner\xi\urcorner,p\rangle \in \text{Proof}$, если $p$ --- гёделев номер
доказательства $\xi$. \pause\\
Обозначим $\pi(x)\equiv\exists p.\psi(x,p)$
\end{dfn} \pause
\begin{dfn}Формулой Consis назовём формулу
$\neg \pi(\overline{\ulcorner 1=0 \urcorner})$
\end{dfn} \pause
Неформальный смысл: <<формальная арифметика непротиворечива>>
\end{frame}

\begin{frame}{Вторая теорема Гёделя о неполноте арифметики}
\begin{thm}Если Consis доказуем, то формальная арифметика противоречива.\end{thm}
\begin{proof}(неформально) \pause
Формулировка 1 теоремы Гёделя о неполноте арифметики: 
<<если Ф.А. непротиворечива, то недоказуемо $\sigma(\overline{\ulcorner\sigma\urcorner})$>>. \pause
То есть, $\forall p.\neg\omega_1(\overline{\ulcorner\sigma\urcorner},p)$. \pause То есть, 
если $\text{Consis}$, то $\sigma(\overline{\ulcorner\sigma\urcorner})$. \pause
То есть, если $\text{Consis}$, то $\sigma(\overline{\ulcorner\sigma\urcorner})$, --- и это можно доказать,
то есть $\vdash\text{Consis}\rightarrow\sigma(\overline{\ulcorner\sigma\urcorner})$. \pause Однако
если формальная арифметика непротиворечива, то $\not\vdash\sigma(\overline{\ulcorner\sigma\urcorner})$.
\end{proof}
\end{frame}

\begin{frame}{Слишком много неформальности}
Рассмотрим такой особый $\text{Consis}'$:
$$\begin{array}{l}\pi'(x) := \exists p.\psi(x,p) \with \neg\psi(\overline{\ulcorner 1 = 0 \urcorner},p)\\
                  \text{Consis}' := \pi'(\overline{\ulcorner 1 = 0 \urcorner})\end{array}$$

Заметим:
\begin{enumerate}
\item Если ФА непротиворечива, то $\llbracket \pi'(x) \rrbracket = \llbracket \pi(x) \rrbracket$:
\begin{itemize}
\item если $x \ne \ulcorner 1 = 0 \urcorner$ и $\llbracket\psi(x,p)\rrbracket = \text{И}$, 
то $\llbracket\psi(\overline{\ulcorner 1 = 0\urcorner},p)\rrbracket = \text{Л}$
\item если $x = \ulcorner 1 = 0 \urcorner$, то $\psi(\overline{\ulcorner 1 = 0 \urcorner},p) = \text{Л}$ при любом $p$.
\end{itemize}
\item Но $\vdash \text{Consis}'$.
\end{enumerate}
\end{frame}

\begin{frame}{Условия выводимости Гильберта-Бернайса-Лёба}
\begin{dfn}
Будем говорить, что формула $\psi$, выражающая отношение Proof, 
формула $\pi$ и формула Consis соответствуют
условиям Гильберта-Бернайса-Лёба, если следующие условия выполнены для любой формулы $\alpha$:

\begin{enumerate}
\item $\vdash \alpha$ влечет $\vdash \pi(\overline{\ulcorner\alpha\urcorner})$
\item $\vdash \pi (\overline{\ulcorner\alpha\urcorner}) \rightarrow \pi(\overline{\ulcorner\pi(\overline{\ulcorner\alpha\urcorner})\urcorner})$
\item $\vdash \pi (\overline{\ulcorner\alpha\rightarrow \beta\urcorner}) \rightarrow \pi(\overline{\ulcorner\alpha\urcorner}) \rightarrow \pi(\overline{\ulcorner\beta\urcorner})$
\end{enumerate}
\end{dfn}
\end{frame}

\begin{frame}{Первая теорема Гёделя о неполноте ещё раз}
\begin{lmm}Лемма об автоссылках. Для любой формулы $\phi(x_1)$ можно построить 
такую замкнутую формулу $\alpha$ (не использующую неаксиоматических предикатных
и функциональных символов), что $\vdash \phi(\overline{\ulcorner\alpha\urcorner}) \leftrightarrow \alpha$.
\end{lmm}

\begin{thm}Существует такая замкнутая формула $\gamma$, что если Ф.А. непротиворечива, то 
$\not\vdash \gamma$, а если Ф.А. $\omega$-непротиворечива, то и $\not\vdash\neg\gamma$.
\end{thm}
\begin{proof}Рассмотрим $\phi(x_1) \equiv \neg\pi(x_1)$. Тогда по лемме об автоссылках существует
$\gamma$, что $\vdash \gamma \leftrightarrow \neg\pi(\overline{\ulcorner\gamma\urcorner})$.

\begin{itemize}
\item Предположим, что $\vdash \gamma$. Тогда $\vdash \gamma \rightarrow \neg\pi(\overline{\ulcorner\gamma\urcorner})$, то есть $\not\vdash\gamma$
\item Предположим, что $\vdash \neg\gamma$. Тогда $\vdash \pi(\overline{\ulcorner\gamma\urcorner})$,
то есть $\vdash \exists p.\psi(\overline{\ulcorner\gamma\urcorner},p)$. Тогда по $\omega$-непротиворечивости
найдётся $p$, что $\vdash \psi(\overline{\ulcorner\gamma\urcorner},\overline{p})$, то есть $\vdash \gamma$.
\end{itemize}
\end{proof}
\end{frame}

\begin{frame}{Доказательство второй теоремы Гёделя}
\begin{enumerate}
\item Пусть $\gamma$ таково, что $\vdash \gamma \leftrightarrow \neg\pi(\overline{\ulcorner\gamma\urcorner})$.
\item Покажем $\pi(\overline{\ulcorner\gamma\urcorner})\vdash \pi(\overline{\ulcorner 1=0\urcorner})$. 

\begin{enumerate}
\item По условию 2, $\vdash \pi(\overline{\ulcorner\gamma\urcorner}) \rightarrow \pi(\overline{\ulcorner\pi(\overline{\ulcorner\gamma\urcorner})\urcorner})$.
По теореме о дедукции $\pi(\overline{\ulcorner\gamma\urcorner})\vdash \pi(\overline{\ulcorner\pi(\overline{\ulcorner\gamma\urcorner})\urcorner})$;

\item Так как $\vdash \pi(\overline{\ulcorner\gamma\urcorner})\rightarrow\neg\gamma$, то 
по условию 1 $\vdash \pi(\overline{\ulcorner\pi(\overline{\ulcorner\gamma\urcorner})\rightarrow\neg\gamma\urcorner})$;

\item По условию 3, $\pi(\overline{\ulcorner\gamma\urcorner})\vdash \pi(\overline{\ulcorner\pi(\overline{\ulcorner\gamma\urcorner})\rightarrow\neg\gamma\urcorner})\rightarrow
\pi(\overline{\ulcorner\pi(\overline{\ulcorner\gamma\urcorner})\urcorner}) \rightarrow 
\pi(\overline{\ulcorner\neg\gamma\urcorner})$;

\item Таким образом, $\pi(\overline{\ulcorner\gamma\urcorner})\vdash\pi(\overline{\ulcorner\neg\gamma\urcorner})$;

\item Однако $\vdash \gamma\rightarrow\neg\gamma\rightarrow 1=0$. Условие 3 (применить два раза) даст $\pi(\overline{\ulcorner\gamma\urcorner})\vdash \pi(\overline{\ulcorner 1=0 \urcorner})$.
\end{enumerate}

\item $\vdash \neg\pi(\overline{\ulcorner 1=0 \urcorner})\rightarrow\neg\pi(\overline{\ulcorner\gamma\urcorner})$ (т. о дедукции, контрапозиция).
\item $\vdash \neg\pi(\overline{\ulcorner 1=0 \urcorner})\rightarrow\gamma$ (определение $\gamma$).
\end{enumerate}
\end{frame}

\begin{frame}{Расширение на другие теории}
\begin{dfn}Теория $\mathcal{S}$ --- расширение теории $\mathcal{T}$, если
из $\vdash_\mathcal{T} \alpha$ следует $\vdash_\mathcal{S} \alpha$\end{dfn}

\begin{dfn}Теория $\mathcal{S}$ --- рекурсивно-аксиоматизируемая, если найдётся теория $\mathcal{S'}$ с тем же языком, что:
\begin{enumerate}
\item $\vdash_\mathcal{S} \alpha$ тогда и только тогда, когда $\vdash_\mathcal{S'} \alpha$;
\item Множество аксиом теории $\mathcal{S'}$ рекурсивно.
\end{enumerate}
\end{dfn}

\begin{thm}Если $\mathcal{S}$ --- непротиворечивое рекурсивно-аксиоматизируемое расширение формальной арифметики, то
в ней можно доказать аналоги теорем Гёделя о неполноте арифметики.
\end{thm}
\end{frame}

\begin{frame}{Сужение: система Робинсона}
\begin{dfn}Теория первого порядка, использующая нелогические функциональные символы $0$, $(+)$ и $(\cdot)$, нелогический
предикатный символ $(=)$ и следующие нелогические аксиомы, называется системой Робинсона.

\vspace{-0.4cm}
$$\begin{array}{ll}
a = a & a = b \rightarrow b = a \\
a = b \rightarrow b = c \rightarrow a = c & a = b \rightarrow a' = b' \\
a' = b' \rightarrow a = b & \neg 0 = a' \\
a = b \rightarrow a + c = b + c \with c + a = c + b & a = b \rightarrow a \cdot c = b \cdot c \with c \cdot a = c \cdot b \\
\neg a = 0 \rightarrow \exists b. a = b' & a + 0 = a\\
a + b' = (a + b)' & a \cdot 0 = 0 \\
a \cdot b' = a \cdot b + a 
\end{array}$$
\end{dfn}

\vspace{-0.3cm}
Система Робинсона неполна: аксиомы --- в точности утверждения, необходимые для доказательства теорем Гёделя.
Система Робинсона не имеет схем аксиом.
\end{frame}

\begin{frame}{Арифметика Пресбургера}
\begin{dfn}Теория первого порядка, использующая нелогические функциональные символы $0$, $1$, $(+)$, нелогический
предикатный символ $(=)$ и следующие нелогические аксиомы, называется арифметикой Пресбургера.

$$\begin{array}{l}
\neg (0 = x + 1) \\
x + 1 = y + 1 \rightarrow x = y\\
x + 0 = x \\
x + (y + 1) = (x + y) + 1\\
(\varphi(0) \with \forall x.\varphi(x) \rightarrow \varphi(x+1)) \rightarrow \forall y.\varphi(y)
\end{array}$$\end{dfn}

\begin{thm}Арифметика Пресбургера разрешима и синтаксически и семантически полна.\end{thm}
\end{frame}

\begin{frame}{Невыразимость доказуемости}
\begin{dfn}
$\text{Д}_\mathcal{S} = \{ \ulcorner\alpha\urcorner\ |\ \vdash_\mathcal{S}\alpha \}$; 
$\text{И}_\mathcal{S} = \{ \ulcorner\alpha\urcorner\ |\ \llbracket\alpha\rrbracket_\mathcal{S} = \text{И} \}$
\end{dfn}

\begin{lmm}Пусть $D(\ulcorner\alpha\urcorner) = \ulcorner\alpha(\overline{\ulcorner\alpha\urcorner})\urcorner$ для
любой формулы $\alpha(x)$. Тогда $D$ представима в формальной арифметике.
\end{lmm}


\begin{thm}Если расширение Ф.А. $\mathcal{S}$ непротиворечиво и $D$ представима в нём, то $\text{Д}_\mathcal{S}$ невыразимо в $\mathcal{S}$\end{thm}

\begin{proof}Пусть $\delta(a,p)$ представляет $D$, и пусть $\sigma(x)$ выражает множество $\text{Д}_\mathcal{S}$ (рассматриваемое как
одноместное отношение).

Пусть $\alpha(x) := \forall p.\delta(x,p)\rightarrow\neg\sigma(p)$. Верно ли, что $\ulcorner\alpha\urcorner\in\text{Д}_\mathcal{S}$?
\end{proof}
\end{frame}


\end{document}
